Here \(\beta=\kappa=1\), and therefore the definitions give
\[
 \delta_{\mathrm{edge},n}=n^{-1/4},
 \qquad
 \delta_{\mathrm{crit},n}
 =n^{-1/[2(1+2+1+1)]}=n^{-1/10}.
\]
If \(\delta_n\gg n^{-1/4}\), Lemma \(\texttt{lem:bandwidth-phases}\) gives
\[
 h_n\asymp(n\delta_n)^{-1/3}.
\]
Moreover,
\[
 \frac{h_n}{\delta_n}
 \asymp(n\delta_n^4)^{-1/3}\longrightarrow0.
\]

We now make the localized Bernoulli bump construction invoked in Theorem \(\texttt{thm:minimax-risk}\) explicit.  Let \(K(u)=(1-|u|)_+\), choose a fixed \(0<\eta\leq\min\{L,1/4\}\), and set
\[
 g_n(a)=\eta h_nK\left(\frac{a-\delta_n}{h_n}\right).
\]
Let \(P_0\) and \(P_1\) have \(X=1\) almost surely, the common conditional treatment density \(\pi(a)=2a\), and conditional outcome laws
\[
 Y\mid A=a\sim\operatorname{Ber}(1/2)
 \quad\text{under }P_0,
 \qquad
 Y\mid A=a\sim\operatorname{Ber}(1/2+g_n(a))
 \quad\text{under }P_1.
\]
For all sufficiently large \(n\), the support \([\delta_n-h_n,\delta_n+h_n]\) lies in \([0,1]\): the lower endpoint is positive because \(h_n/\delta_n\to0\), and the upper endpoint is below one because \(\delta_n\leq\bar\delta<1\) and \(h_n\to0\).  Both regressions take values in \([0,1]\).  The function \(g_n\) is \(\eta\)-Lipschitz, so the two regressions satisfy Assumption \(\texttt{ass:holder-regression}\) for \(\beta=1\) with radius \(L\).  The density \(\pi(a)=2a\) integrates to one, satisfies Assumption \(\texttt{ass:conditional-density-law}\), and satisfies Assumption \(\texttt{ass:polynomial-thinning}\) for \(\kappa=1\), because the declared ranges give \(c_-\leq2\leq c_+\).  The one-cell mass is one and hence satisfies Assumption \(\texttt{ass:stratum-mass}\).  Thus Definition \(\texttt{def:model-class}\) shows that \(P_0,P_1\in\mathcal M\); in particular, these are same-class alternatives.

For \(0\leq g\leq1/4\), define
\[
 f(g)=\operatorname{KL}\{\operatorname{Ber}(1/2+g),\operatorname{Ber}(1/2)\}.
\]
Direct differentiation gives \(f(0)=f'(0)=0\) and
\[
 f''(g)=\frac{1}{(1/2+g)(1/2-g)}.
\]
This second derivative is bounded above and away from zero on \([0,1/4]\).  The identity
\[
 f(g)=\int_0^g(g-t)f''(t)\,\mathrm dt
\]
therefore gives constants \(0<c<C<\infty\), independent of \(g\), such that
\[
 c g^2\leq
 \operatorname{KL}\{\operatorname{Ber}(1/2+g),\operatorname{Ber}(1/2)\}
 \leq Cg^2.
\]
The two laws have the same treatment margin.  Since \(a\asymp\delta_n\) uniformly on the support of \(g_n\), conditioning on \(A\) and changing variables \(u=(a-\delta_n)/h_n\) show that the one-observation divergence satisfies
\[
 \begin{aligned}
 \operatorname{KL}(P_1,P_0)
 &\asymp \int_0^1 2a\,g_n(a)^2\,\mathrm da \\
 &\asymp \delta_n h_n^3
 \int_{-1}^1K(u)^2\,\mathrm du
 \asymp\delta_nh_n^3.
 \end{aligned}
\]
Assumption \(\texttt{ass:iid-sampling}\) makes the product log-likelihood ratio the sum of its \(n\) one-observation counterparts.  Taking expectation, equivalently using the product-divergence calculation in the localized lower bound of Theorem \(\texttt{thm:minimax-risk}\), therefore gives
\[
 \operatorname{KL}(P_1^{\otimes n},P_0^{\otimes n})
 =n\operatorname{KL}(P_1,P_0)
 \asymp n\delta_nh_n^3.
\]
The bandwidth relation makes this quantity of constant order.

For the target separation, the definition of \(\theta_{\delta_n}\), together with \(\pi(a)=2a\), gives
\[
 q(\delta_n)=\int_0^{\delta_n}2a\,\mathrm da=\delta_n^2
\]
and
\[
 \begin{aligned}
 \theta_{\delta_n}(P_1)-\theta_{\delta_n}(P_0)
 &=q(\delta_n)g_n(\delta_n)
   +\int_{\delta_n}^1 2a\,g_n(a)\,\mathrm da \\
 &=\eta\delta_n^2h_n+O(\delta_nh_n^2).
 \end{aligned}
\]
Both displayed contributions are nonnegative, and \(h_n/\delta_n\to0\); hence the separation is \(\asymp\delta_n^2h_n\).  This is the explicit localized bump construction used by the minimax lower bound, so the reference is unambiguous.

Finally, equating the atom separation with the root-\(n\) term and substituting the bandwidth order gives
\[
 \delta_n^2(n\delta_n)^{-1/3}=n^{-1/2}
 \quad\Longleftrightarrow\quad
 \delta_n^{5/3}=n^{-1/6}
 \quad\Longleftrightarrow\quad
 \delta_n=n^{-1/10}.
\]
This is exactly the transition asserted in Theorem \(\texttt{thm:phase-diagram}\).